\subsection{Experimental Evaluation and Offset Calibration}
\label{app:jaxka-experimental-evaluation}

The reported experimental benchmark used PKAD-3 entries distributed with the twenty KaML train/test splits.\supercite{shenKaMLsPredictingProtein2025} We retained entries with a prepared RCSB chain, excluding 103 entries without such a structure, including entries identified only by mutant labels. This left 1,604 entries. The common-entry comparison required predictions from each of the benchmark's GQT, PypKa, pKAI, pKAI+, PROPKA and null arms, leaving 1,421 entries. KaML's split-specific predictions were not part of this pooled comparison. All-atom predictors were evaluated on deposited and FASPR-repacked side chains separately. GQT used backbone coordinates and sequence. The null assigned the model-compound value recorded with each KaML entry.

Pooled errors counted each retained benchmark entry once. Shift-stratified errors used each predictor's available entries and grouped them by the absolute difference between experimental and model-compound $\mathrm{p}K_a$, with boundaries at one and two units. Run4 covered 1,489 entries. These partitions therefore used a different cohort from the common-entry comparison. Neither pooled analysis fitted calibration parameters.

Calibration tested whether residue-type offsets explained part of the experimental error without changing the structural predictor. For each split and residue type, we calculated the mean prediction-minus-measurement residual over that split's training entries with available predictions, then subtracted it from predictions for the corresponding test entries. An absent training residue type received zero offset. Test measurements did not enter offset fitting. We calculated test RMSE within each split, then reported the mean and population standard deviation over twenty splits, separately for all sites, acidic groups (Asp, Glu, Cys and Tyr) and basic groups (His and Lys). Coverage is the fraction of test-entry occurrences with predictions across those splits. These repeated partitions are not twenty independent experimental datasets.

The experimental measurements did not supervise GQT training. The benchmark record nevertheless allows overlap between evaluated structures and the theoretical-label training data used by GQT, pKAI or pKAI+. It therefore tests transfer from computational labels to measurements without establishing a structure-disjoint comparison across all predictors. The pooled and calibrated results are given in \cref{tab:jaxka-experimental-common,tab:jaxka-experimental-diagnostics}.
